%!TeX TS-program = lualatex
%!TeX encoding = UTF-8 Unicode
%!TeX spellcheck = en-UK
% -*- coding: UTF-8; -*-
% vim: set fenc=utf-8
%=== file: myfigure.tex ==========================================
\documentclass[tikz,border=2pt]{standalone}   % “border” adds optional margin
% -------------------------------------------------------------
% 1. Packages you used in the original doc (add any that are missing)
\usepackage{amsmath,amssymb}          % math, symbols
\usepackage{tikz}
\usetikzlibrary{arrows,calc,positioning,shapes,decorations.pathreplacing} % etc.
\usepackage{pgfplots}                % if you used pgfplots
\pgfplotsset{compat=newest}           % optional – keep the same version as before
% -------------------------------------------------------------
% 2. Any global TikZ settings that the picture depends on
\tikzset{
   % example: thin arrows everywhere
   > = stealth,
   % ... copy from your original preamble
}
% -------------------------------------------------------------
\begin{document}
% ------------------- START OF YOUR PICTURE --------------------
  \begin{tikzpicture}[
    x=1cm,
    y=1cm,
    line cap=round,
    line join=round,
    font=\sffamily,
    rect/.style={draw=#1, fill=#1!10, line width=0.9pt, rounded corners=2.5pt, minimum width=2.1cm, minimum height=0.9cm, inner sep=1pt},
    spike/.style={draw=black, line width=1.05pt, -{stealth[length=2.8mm,width=2.2mm]}},
    beta/.style={draw=cyan!60!black, line width=1.2pt, -{stealth[length=2.5mm,width=2.0mm]}},
    weights/.style={draw=blue!80!black, line width=1.3pt, -{stealth[length=2.6mm,width=2.0mm]}},
    thr/.style={draw=orange!55!brown!85!black, line width=1.3pt, -{stealth[length=2.6mm,width=2.0mm]}},
    labelbeta/.style={above, font=\large, text=cyan!60!black, midway},
    labelweights/.style={above, font=\large, text=blue!80!black, pos=0.50},
    labelthr/.style={above, sloped, font=\large, text=orange!55!brown!85!black, midway},
  ]
 
  \def\colSep{3.55}
  \def\xY{-2.05}
  \def\iY{-0.55}
  \def\uY{0.95}
  \def\sY{2.45}
  \newcommand{\colpos}[2]{({\colSep*#1,#2})}
  \def\tLabelY{-3.0}
  \def\ellipsisX{7.95}
  \def\blueRise{0.95}
  \def\blueRun{1.15}
 
  % Positions for the three time steps
  % (Removed manual coordinate definitions)
 
  % Dotted recurrent shell over U and S layers
  \path[draw=black, dashed, line width=1.15pt, rounded corners=0.58cm, fill=black!4, fill opacity=0.20] (-1.15,-0.05) rectangle (7.75,3.15);
 
  % Layer nodes
  \foreach \i in {0,1,2} {
    \node[rect=blue!55!black] at (\colSep*\i, \xY) {X[\i]};
    \node[rect=cyan!60!black] at (\colSep*\i, \iY) {$I_{\mathrm{in}}[\i]$};
    \node[rect=orange!70!brown!85!black] at (\colSep*\i, \sY) {S[\i]};
    \node[rect=red!55!black] at (\colSep*\i, \uY) {U[\i]};
  }
 
  % Vertical input chain
  \foreach \i in {0,1,2} {
    \draw[spike] (\colSep*\i, \xY) -- node[midway, right=2pt, font=\large] {$W_0$} (\colSep*\i, \iY);
    \draw[spike] (\colSep*\i, \iY) -- (\colSep*\i, \uY);
    \draw[spike] (\colSep*\i, \uY) -- (\colSep*\i, \sY);
    \draw[spike] (\colSep*\i, \sY) -- ++(0,0.6);
  }
 
  % Recurrent beta and threshold branches
  \foreach \i [evaluate=\i as \j using int(\i+1)] in {0,1} {
    \draw[beta] (\colSep*\i, \uY) -- node[labelbeta] {$\beta$} (\colSep*\j, \uY);
    \draw[thr] (\colSep*\i, \sY) -- node[labelthr] {$-U_{\mathrm{thr}}$} (\colSep*\j, \uY);
  }
 
  % Input weights between successive columns
  \draw[weights] (-1.55,-1.25) -- node[labelweights] {$W_{-1}$} (0, \iY);
  \draw[weights, shorten <=1.5pt, shorten >=1.5pt] (\colSep*0, \xY) -- node[labelweights] {$W_{-1}$} (\colSep*1, \iY);
  \draw[weights, shorten <=1.5pt, shorten >=1.5pt] (\colSep*1, \xY) -- node[labelweights] {$W_{-1}$} (\colSep*2, \iY);
  \draw[weights] (\colSep*2, \xY) -- ++(\blueRun,\blueRise);
 
  % Time labels
  \node[font=\Large] at (0,\tLabelY) {$t = 0$};
  \node[font=\Large] at (\colSep,\tLabelY) {$t = 1$};
  \node[font=\Large] at (2*\colSep,\tLabelY) {$t = 2$};
  \node[font=\Large] at (\ellipsisX,0.0) {$\cdots$};
  \node[font=\Large] at (\ellipsisX,1.0) {$\cdots$};
 
  % Top arrows and labels
  \node[font=\Large] at (\ellipsisX,2.85) {$\cdots$};
 
  % Clip-friendly top border for the dashed envelope
  \path[use as bounding box] (-1.7,-3.35) rectangle (9.8,3.55);
  \end{tikzpicture}%
% -------------------- END OF YOUR PICTURE -----------------
\end{document}
%================================================================
